\documentclass{article}

\title{Ch. 7.3 HW}
\author{Principles of Microeconomics}
\date{\today}

\begin{document}

\maketitle

\subsection*{Problems}
Quick Quiz \#7-10, Problems \& Applications \#6-8

\subsection*{Quick Quiz}

\begin{enumerate}
\setcounter{enumi}{6}

\item Isabelle values her time at \$60 an hour. She spends 2 hours giving Jayla a massage. Jayla was willing to pay as much as \$300 for the massage, but they negotiated a price of \$200. In this transaction,

	\begin{enumerate}
	
	\item consumer surplus is \$20 larger than producer surplus.
	
	\item consumer surplus is \$40 larger than producer surplus.
	
	\item producer surplus is \$20 larger than consumer surplus.
	
	\item producer surplus is \$40 larger than consumer surplus.
	
	\end{enumerate}
	
\item An efficient allocation of resources maximizes

	\begin{enumerate}
	
	\item consumer surplus.
	
	\item producer surplus.
	
	\item consumer surplus plus producer surplus.
	
	\item consumer surplus minus producer surplus.
	
	\end{enumerate}
	
\item When a market is in equilibrium, the buyers are those with the \underline{\phantom{xxxxxxxxxx}} willingness to pay, and the sellers are those with the \underline{\phantom{xxxxxxxxxx}} costs.

	\begin{enumerate}
	
	\item highest; highest
	
	\item highest; lowest
	
	\item lowest; highest
	
	\item lowest; lowest
	
	\end{enumerate}
	
\newpage
	
\item Producing a quantity larger than the equilibrium of supply and demand is inefficient because the marginal buyer's willingness to pay is 

	\begin{enumerate}
	
	\item negative.
	
	\item zero.
	
	\item positive but less than the marginal seller's cost.
	
	\item positive and greater than the marginal seller's cost.
	
	\end{enumerate}

\end{enumerate}

\subsection*{Problems and Applications}

\begin{enumerate}
\setcounter{enumi}{5}

\item Consider a market in which Bert from problem 4 is the buyer and Ernie from problem 5 is the seller.

	\begin{enumerate}
	
	\item Use Ernie's supply schedule and Bert's demand schedule to find the quantity supplied and quantity demanded at prices of \$2, \$4, and \$6. Which of these prices brings supply and demand into equilibrium?
	
	\item What are consumer surplus, producer surplus, and total surplus in this equilibrium?
	
	\item If Ernie produced and Bert consumed one fewer bottle of water, what would happen to total surplus?
	
	\item If Ernie produced and Bert consumed one additional bottle of water, what would happen to total surplus?
	
	\end{enumerate}
	
\item The cost of producing flat-screen TVs has fallen over the past decade. Let's consider some implications of this change.

	\begin{enumerate}
	
	\item Draw a supply-and-demand diagram to show the effect of falling production costs on the price and quantity of flat-screen TVs sold.
	
	\item In your diagram, show what happens to consumer surplus and producer surplus.
	
	\item Suppose the supply of flat screen TVs is very elastic. Who benefits most from falling production costs -- consumers or producers of these TVs?
	
	\end{enumerate}
	
\item Four consumers are willing to pay these amounts for haircuts:

	\begin{center}
	\begin{tabular}{| c | c | c | c |}
	\hline
	Gloria: \$35 & Jay: \$10 & Claire: \$40 & Phil: \$25 \\
	\hline
	\end{tabular}
	\end{center}
	
	Four haircutting businesses have these costs:
	
	\begin{center}
	\begin{tabular}{| c | c | c | c |}
	\hline
	Firm A: \$15 & Firm B: \$30 & Firm C: \$20 & Firm D: \$10 \\
	\hline
	\end{tabular}
	\end{center}
	
	Each firm can give, at most, one haircut. To achieve efficiency, how many haircuts should be given? Which businesses should cut hair and which consumers should have their hair cut? How large is the maximum possible total surplus?

\end{enumerate}

\end{document}